\documentclass{article} % For LaTeX2e
\usepackage{nips15submit_e,times}
\usepackage{hyperref}
\usepackage{url}
\usepackage{amsmath,amssymb,graphicx,booktabs}
\usepackage{algorithm}
\usepackage{algpseudocode}

% Course-required: use NIPS camera-ready format.
\nipsfinalcopy

\title{Quantile-Initialised Ordinal Collaborative Filtering\\ with Gated Auxiliary Fusion}
\author{
  Kevin Zhu \\ Australian National University \\ \texttt{u7338066@anu.edu.au}
}

\begin{document}
\maketitle

\begin{abstract}
We propose a simple calibration recipe for cumulative-link collaborative
filtering: initialise the monotone thresholds from the empirical rating
quantiles of the training set, parameterise their spacing through a
softplus-inverse, and exclude them from $L_2$ regularisation. The recipe
matches the data marginal exactly at the start of training and, on a
transparent dot-product matrix-factorisation backbone, achieves low
test negative log-likelihood ($1.2584$ across the five canonical
MovieLens-100K splits, vs the $\log 5{\approx}1.61$ uniform baseline)
without sacrificing RMSE. Paired with per-dimension gated fusion of side
information, the gated$+$ordinal configuration attains
RMSE $0.9124\pm0.0036$ across the five splits, beating both a no-fusion MF
analogue ($0.9160$) and additive fusion ($0.9189$) on the across-split mean.
The contribution is practical rather than algorithmic: standard components
are combined into a reproducible, calibrated recommender pipeline.
\end{abstract}

\section{Introduction}
\label{sec:intro}

Collaborative filtering (CF) for explicit feedback is commonly framed as
predicting an unobserved scalar rating $r_{ui}$ from a sparse user--item rating
matrix $R\in\mathbb{R}^{n\times m}$. Latent-factor models based on matrix
factorisation~\cite{koren2009mf} remain standard baselines and building blocks
on MovieLens~\cite{harper2015movielens}.
Two persistent issues remain when these models are applied as a course-style
data-analytics tool rather than a production recommender. First, side
information (user demographics, item genres) is typically injected by adding a
projected feature vector to the latent embedding, with no learned mechanism to
decide \emph{when} to trust the embedding versus the side information. Second,
the rating is treated as a continuous quantity even though it is in fact a
five-class ordinal variable; the resulting predictions can be RMSE-competitive
yet badly calibrated as distributions over star ratings.

\paragraph{Related work.}
Gated content fusion has been explored by GATE~\cite{ma2019gate}; ordinal
regression starts from McCullagh's cumulative-link model~\cite{mccullagh1980ordinal}
and has been applied to MovieLens by OrdRec~\cite{koren2011ordrec}. Recent
ordinal recommender baselines include OrdNMF~\cite{gouvert2020ordnmf} and
OPRFM~\cite{zaman2025oprfm}. Our contribution is not a new family of
models, but a compact integration of gated side information and calibrated
cumulative-link prediction on a transparent MF core.

\paragraph{Contributions.}
We make two practical contributions. First, we initialise cumulative-link
thresholds from the training-set rating quantiles and exclude those thresholds
from weight decay, so the ordinal head starts with the correct rating
marginal. Second, we pair that head with symmetric per-dimension gates that
blend ID embeddings with user and item side-feature projections. The ablation
separates their roles: gating gives the larger RMSE change, while the ordinal
head gives calibrated class probabilities and better boundary-class behaviour.

\section{Problem Definition}
\label{sec:problem}

Let $\mathcal{U}=\{1,\dots,n\}$ be the set of users and $\mathcal{I}=\{1,\dots,m\}$
the set of items. Observed ratings form a sparse set
$\mathcal{D}=\{(u_t,i_t,r_t)\}_{t=1}^{N}$ with $r_t\in\{1,2,3,4,5\}$. Each user
$u$ carries a side-information vector $f_u\in\mathbb{R}^{d_u}$ (z-scored age,
one-hot gender, one-hot occupation in MovieLens-100K, $d_u=24$) and each
item $i$ carries $g_i\in\mathbb{R}^{d_i}$ (multi-hot genre indicator,
$d_i=19$). The task
is to predict $\hat r_{ui}$ for held-out $(u,i)$ pairs, evaluated by RMSE, MAE,
classification accuracy of the rounded prediction, and where applicable the
held-out negative log-likelihood (NLL) of the predicted rating distribution.

\section{Method}
\label{sec:method}

\subsection{Matrix factorisation backbone}
We start from a biased dot-product MF model. Each user and item carries a
$d$-dimensional embedding ($p_u, q_i\in\mathbb{R}^d$) and a scalar bias
($b_u, b_i\in\mathbb{R}$). The pre-head score is
\begin{equation}
\label{eq:score}
s_{ui} \;=\; p'_u \cdot q'_i \;+\; b_u \;+\; b_i,
\end{equation}
where $p'_u,q'_i$ are the (possibly fused) representations defined below. With
fusion disabled, $p'_u\equiv p_u$ and $q'_i\equiv q_i$.

\subsection{Gated auxiliary fusion}
\label{sec:method:fusion}
Let $a_u=\mathrm{ReLU}(W_u f_u + b_{f,u})$ be a learned projection of the user
side features into the embedding space, and similarly $a_i=\mathrm{ReLU}(W_i g_i + b_{f,i})$
for items. We define the fused user representation as
\begin{equation}
\label{eq:gate}
g_u \;=\; \sigma\!\left(W_g \,[\,p_u \,\Vert\, a_u\,] + b_g\right),\qquad
p'_u \;=\; (\mathbf{1}-g_u)\odot p_u \;+\; g_u\odot a_u,
\end{equation}
where $\Vert$ denotes concatenation and $\odot$ is the Hadamard product, so the
gate $g_u\in[0,1]^d$ acts \emph{per dimension}. The item-side fusion is
identical with $(q_i,a_i)$ in place of $(p_u,a_u)$. We zero-initialise both
$W_g$ and $b_g$, which makes $g_u\equiv \tfrac12\mathbf{1}$ at training start.
This gives every dimension the same neutral starting mixture between the ID
embedding and the ReLU-projected side features; the gate then opens and closes
per dimension as gradient flows. The ablation compares this against
\emph{additive} fusion ($p'_u = p_u + a_u$) and \emph{none} fusion
($p'_u\equiv p_u$), with analogous item-side definitions.

\subsection{Ordinal regression head}
\label{sec:method:head}
For the head we treat $r$ as an ordinal variable with $K=5$ classes. We use the
cumulative-link model of McCullagh~\cite{mccullagh1980ordinal} with monotone
thresholds $\theta_1<\theta_2<\theta_3<\theta_4$:
\begin{equation}
\label{eq:cum}
P(r\le k\mid s) = \sigma(\theta_k - s),\qquad
P(r=k\mid s) = \sigma(\theta_k - s)-\sigma(\theta_{k-1} - s),
\end{equation}
with $\theta_0=-\infty$ and $\theta_5=+\infty$. We enforce monotonicity by the
softplus reparameterisation
\begin{equation}
\label{eq:softplus}
\theta_k=\theta_1+\sum_{j=1}^{k-1}\mathrm{softplus}(\delta_j),
\end{equation}
which guarantees $\theta_k<\theta_{k+1}$ for any unconstrained parameters.

\paragraph{Empirical-quantile initialisation.}
We initialise from the empirical
rating distribution. Let $\hat p_k$ be the training-set frequency of rating
$k$ and $\hat F_k=\sum_{j\le k}\hat p_j$ its cumulative. Choosing
\begin{equation}
\label{eq:quantile-init}
\theta_k = \mathrm{logit}(\hat F_k),\qquad
\delta_j = \log\!\bigl(\exp(\theta_{j+1}-\theta_j) - 1\bigr)
\end{equation}
guarantees $P(r{\le}k\mid s{=}0)=\hat F_k$, so the predicted marginal at the
start of training equals the data marginal exactly. The second equation is
just the softplus inverse, recovering the unconstrained $\delta_j$ that the
optimiser then updates. We also exclude $\{\theta_1,\delta_j\}$ from weight
decay because their scale carries calibration information.

The training loss is the negative log-likelihood
$\mathcal{L}_{\mathrm{ord}}=-\frac{1}{N}\sum_t \log P(r=r_t\mid s_{u_ti_t})$,
and the point prediction is the expected rating
$\hat r_{ui}=\sum_{k=1}^{5}k\,P(r=k\mid s_{ui})$. As an ablation we also
consider the standard \emph{sigmoid} head:
$\hat r_{ui}=4\sigma(s_{ui})+1$ trained with mean-squared error.

\begin{algorithm}[ht]
\caption{Forward pass + ordinal NLL for one observed rating $(u,i,r)$.}
\label{alg:forward}
\begin{algorithmic}[1]
\State $a_u\!\gets\!\mathrm{ReLU}(W_u f_u {+}b_{f,u})$;\;
       $a_i\!\gets\!\mathrm{ReLU}(W_i g_i {+}b_{f,i})$
       \Comment{side projections}
\State $g_u\!\gets\!\sigma(W_g^{(u)}[E_u[u]\Vert a_u]{+}b_g^{(u)})$;\;
       $g_i\!\gets\!\sigma(W_g^{(i)}[E_i[i]\Vert a_i]{+}b_g^{(i)})$
       \Comment{per-dim gates}
\State $p'\!\gets\!(\mathbf{1}{-}g_u)\!\odot\!E_u[u]+g_u\!\odot\! a_u$;\;
       $q'\!\gets\!(\mathbf{1}{-}g_i)\!\odot\!E_i[i]+g_i\!\odot\! a_i$
\State $s\!\gets\!p'\!\cdot\!q'+b_u+b_i$;\;
       $\theta_k\!\gets\!\theta_1+\sum_{j<k}\mathrm{softplus}(\delta_j)$
\State $P(r{=}k\mid s)\!\gets\!\sigma(\theta_k{-}s)-\sigma(\theta_{k-1}{-}s)$,\;
       with $\theta_0{=}-\infty,\theta_K{=}\infty$
\State \Return $\mathcal{L}=-\log P(r{=}r_{\text{true}}\mid s)$
\end{algorithmic}
\end{algorithm}

\subsection{Complexity}
\label{sec:complexity}
The MF backbone has $(n+m)d$ embedding parameters plus $n+m$ biases; for
MovieLens-100K with $n=943$, $m=1682$, $d=128$, that is $\approx 3.4\times10^5$
parameters. Each fusion module adds a projection ($d_{u/i}\cdot d + d$) and a gate
($2d^2+d$); summed over both sides this is $\approx 7\times10^4$ parameters
--- about one-fifth of the embedding budget. The ordinal head adds only
$K-1=4$ parameters, and the dominant per-rating cost remains $O(d)$.

\section{Experiments}
\label{sec:exp}

\subsection{Setup}
We use MovieLens-100K~\cite{harper2015movielens} on its five canonical
$80/20$ splits \texttt{u1}--\texttt{u5}, each yielding $80{,}000$ training and
$20{,}000$ test ratings on the same $943$ users and $1{,}682$ items.
Headline numbers are reported as mean $\pm$ std across the five per-split
means (each split-mean itself averages three training seeds), giving
$15$ trained models per ablation cell. User side features are z-scored age,
two-class one-hot gender, and 21-class one-hot occupation ($d_u=24$); item
features are 19-genre multi-hot ($d_i=19$). All experiments use Adam with
learning rate $10^{-3}$, $d=128$, and weight decay $\lambda=10^{-5}$ on all
parameters except the ordinal-head thresholds. Early stopping is driven by
validation RMSE only, and the test set is read once after restoring the
best-validation weights.

\subsection{Baselines and metrics}
The ablation table includes the no-fusion sigmoid model as a biased-MF
analogue, additive fusion, and gated fusion, each under sigmoid and ordinal
heads. Metrics are RMSE, MAE, rounded-rating accuracy, and NLL where a native
class-probability head is available.

\subsection{Results}
As a sanity check, the proposed model also beats three transparent baselines
on the \texttt{u1} split under the headline protocol (single seed, patience
30), shown in Table~\ref{tab:baselines}. The multi-split ablation below is the
main evidence, but this table guards against only comparing the model to its
own variants.

\begin{table}[ht]
\centering
\caption{Transparent baselines on MovieLens-100K \texttt{u1}. Lower is better
for RMSE/MAE/NLL; higher is better for accuracy.}
\label{tab:baselines}
\begin{tabular}{lcccc}
\toprule
Model & RMSE & MAE & Accuracy & NLL \\
\midrule
SVD (rank 20) & $1.0834$ & $0.8896$ & $0.3716$ & --- \\
Biased MF & $0.9221$ & $0.7274$ & $0.4124$ & --- \\
NMF & $0.9197$ & $0.7252$ & $0.4146$ & --- \\
Gated$+$ordinal & $\mathbf{0.9122}$ & $\mathbf{0.7130}$ & $\mathbf{0.4367}$ & $\mathbf{1.2472}$ \\
\bottomrule
\end{tabular}
\end{table}

The ablation matrix in Table~\ref{tab:ablation} reports the full
$\{$none, additive, gated$\}\times\{$sigmoid, ordinal$\}$ cross-product on
three seeds. The proposed configuration --- gated fusion with the ordinal head --- has the
best mean value on all four reported metrics. Across the five
splits it improves RMSE by $0.0036$ over the no-fusion baseline
(\emph{none}+sigmoid, our analogue of plain MF) and by $0.0065$ over the
additive-fusion variant under the sigmoid head. Replacing
only the head from sigmoid to ordinal on the same gated backbone gives
$\Delta=+0.0011$ RMSE (within between-split std), and improves MAE by
$0.0056$ and accuracy by $0.0087$. The ordinal cell is also the only one
that emits a native class distribution (NLL $1.2584$).
Because the RMSE differences are small relative to between-split variation, we
interpret the main value of the ordinal head as probabilistic output and
per-class behaviour rather than a large RMSE improvement.

\begin{table}[ht]
\centering
\caption{Ablation matrix on MovieLens-100K, mean $\pm$ std across the five
canonical splits \texttt{u1}--\texttt{u5} (each split-mean averages three
seeds). Bold = best per column. Lower is better for RMSE/MAE/NLL, higher for
Accuracy.}
\label{tab:ablation}
\resizebox{\linewidth}{!}{%
\begin{tabular}{llcccc}
\toprule
Fusion & Head & RMSE & MAE & Accuracy & NLL \\
\midrule
none      & sigmoid & $0.9160\pm0.0061$ & $0.7241\pm0.0059$ & $0.4233\pm0.0044$ & --- \\
none      & ordinal & $0.9170\pm0.0072$ & $0.7204\pm0.0069$ & $0.4295\pm0.0056$ & $1.2634\pm0.0050$ \\
additive  & sigmoid & $0.9189\pm0.0044$ & $0.7238\pm0.0055$ & $0.4243\pm0.0057$ & --- \\
additive  & ordinal & $0.9174\pm0.0040$ & $0.7173\pm0.0035$ & $0.4340\pm0.0038$ & $1.2654\pm0.0017$ \\
gated     & sigmoid & $0.9135\pm0.0037$ & $0.7182\pm0.0034$ & $0.4286\pm0.0054$ & --- \\
\textbf{gated} & \textbf{ordinal} & $\mathbf{0.9124\pm0.0036}$ & $\mathbf{0.7126\pm0.0042}$ & $\mathbf{0.4373\pm0.0045}$ & $\mathbf{1.2584\pm0.0043}$ \\
\bottomrule
\end{tabular}%
}
\end{table}

\paragraph{Interpretation.}
Table~\ref{tab:ablation} supports two claims. First, gated fusion is the main
RMSE mechanism: holding the head fixed, replacing \emph{none} with
\emph{gated} drops RMSE by $0.0025$ (sigmoid) or $0.0046$ (ordinal), while
replacing \emph{additive} with \emph{gated} drops it by about $0.005$ under
either head. Second, the ordinal head is mainly a probability/calibration
improvement. Switching sigmoid to ordinal changes RMSE by at most $0.0015$,
but consistently improves MAE and accuracy and is the only head with native
NLL. The proposed configuration's NLL $1.2584$ is well below the
$\log 5\approx1.609$ uniform baseline.

\paragraph{Class-level behaviour.}
On \texttt{u1.test}, the ordinal head improves macro-F1 from $0.341$ to
$0.369$ on the same gated backbone, mostly by recovering boundary ratings
(class-1 recall $0.169$ vs $0.102$, class-5 recall $0.235$ vs $0.204$).
Per-class ECE shows the same pattern: the ordinal head is better on classes
1, 2, 3, and 5, while the kernelised sigmoid head is better on the dominant
class 4. This is why we treat the ordinal head as a probability/class-behaviour
improvement rather than a large RMSE improvement.

\begin{figure}[ht]
\centering
\includegraphics[width=0.62\linewidth]{figures/per_class_f1.png}
\caption{Per-class F1 on \texttt{u1.test}. The ordinal head mainly improves
the boundary classes, rather than creating a uniform gain across all ratings.}
\label{fig:perclass}
\end{figure}

The calibration view is consistent with the class-level result: pooled ECE is
similar for the kernelised sigmoid and ordinal heads, but the ordinal head is
better on classes 1, 2, 3, and 5, while the sigmoid head wins only on the
dominant class 4.

\begin{figure}[ht]
\centering
\includegraphics[width=0.82\linewidth]{figures/reliability_curves.png}
\caption{Reliability diagrams on \texttt{u1.test}. The ordinal head spreads
probability mass more evenly across rating classes, while the sigmoid head is
strongest on the dominant class.}
\label{fig:reliability}
\end{figure}

Figure~\ref{fig:sens} shows the sensitivity sweep over embedding dimension
$d$ and weight decay $\lambda$. The chosen $d=128$ is the elbow: $d=256$
improves RMSE by only $0.0024$ at $34\%$ extra wall time, comparable to
seed-to-seed variance. The best tested regularisation is
$\lambda=10^{-5}$; $\lambda=10^{-4}$ overregularises by $0.0080$ RMSE.

\begin{figure}[ht]
\centering
\includegraphics[width=0.92\linewidth]{figures/sensitivity_curves.png}
\caption{Sensitivity of gated$+$ordinal test RMSE to embedding dimension
$d$ and weight decay $\lambda$ on \texttt{u1}; error bars are $\pm1$ std over
three seeds.}
\label{fig:sens}
\end{figure}

\paragraph{Embedding and regime diagnostics.}
The post-fusion embeddings contain side-information structure without an
auxiliary classification loss. In 128-D, the silhouette score is $+0.0848$ for
user occupation and $+0.1930$ for item dominant genre. After 2-D projection,
item genre remains visible under IsoMap
($+0.232$) but is much weaker under PCA ($+0.048$),
suggesting that the item-side genre geometry is curved. The user-side
occupation silhouette becomes negative under all 2-D projections, so those
features appear distributed across too many dimensions to visualise faithfully
in a plane. This diagnostic is secondary to the RMSE/NLL result, but it checks
that the gate is not merely adding opaque parameters.

\begin{figure}[ht]
\centering
\includegraphics[width=0.80\linewidth]{figures/manifold_grid_white.png}
\caption{Two-dimensional projections of post-fusion user embeddings (top,
occupation) and item embeddings (bottom, dominant genre). Columns show PCA,
classical MDS, and IsoMap; the item-side genre pattern is clearest under the
nonlinear projection.}
\label{fig:manifold}
\end{figure}

\begin{table}[ht]
\centering
\caption{Cold-start RMSE by user activity. Negative $\Delta$ means additive
fusion wins.}
\label{tab:coldstart}
\begin{tabular}{lrrrrr}
\toprule
$|R_u|$ bucket & $<\!30$ & $30$--$59$ & $60$--$119$ & $120$--$239$ & $\ge\!240$ \\
\midrule
gated RMSE & $0.9799$ & $0.9649$ & $0.8841$ & $0.8896$ & $0.9200$ \\
additive RMSE & $0.9803$ & $0.9768$ & $0.8929$ & $0.8946$ & $0.9128$ \\
$\Delta$ & $+0.0004$ & $+0.0119$ & $+0.0088$ & $+0.0051$ & $-0.0071$ \\
mean $g_u$ & $0.246$ & $0.256$ & $0.263$ & $0.268$ & $0.291$ \\
\bottomrule
\end{tabular}
\end{table}

The cold-start breakdown in Table~\ref{tab:coldstart} shows that gated fusion
is not uniformly beneficial: it is strongest for mid-activity users, ties on
the coldest bucket, and reverses for the 14 most active users. A three-cell
MovieLens-1M check preserved the same ordering at larger scale, with
gated$+$ordinal reaching RMSE $0.8481$ and NLL $1.1842$.

\begin{table}[ht]
\centering
\caption{Scale-up to MovieLens-1M, random $80/10/10$ split, mean $\pm$ std over
three seeds.}
\label{tab:ml1m}
\resizebox{\linewidth}{!}{%
\begin{tabular}{llcccc}
\toprule
Fusion & Head & RMSE & MAE & Accuracy & NLL \\
\midrule
none & sigmoid & $0.8567\pm0.0008$ & $0.6782\pm0.0007$ & $0.4480\pm0.0005$ & --- \\
gated & sigmoid & $0.8502\pm0.0003$ & $0.6695\pm0.0001$ & $0.4579\pm0.0002$ & --- \\
\textbf{gated} & \textbf{ordinal} & $\mathbf{0.8481\pm0.0006}$ & $\mathbf{0.6630\pm0.0010}$ & $\mathbf{0.4654\pm0.0014}$ & $\mathbf{1.1842\pm0.0013}$ \\
\bottomrule
\end{tabular}%
}
\end{table}

\subsection{Reproducibility}
The code path is intentionally small enough to audit. All tables and figures in
the report are regenerated from archived JSON/CSV outputs or from the scripts
in Table~\ref{tab:repro}. Each script records the model configuration, split,
seed, and metric outputs used in the paper. The reported ablation uses training
seeds $\{42,43,44\}$ on each canonical split and fixes the validation-split
seed to avoid mixing split variance with model-initialisation variance.

\begin{table}[ht]
\centering
\caption{Core reproduction scripts. Wall times are approximate on M1 Pro CPU.}
\label{tab:repro}
\resizebox{\linewidth}{!}{%
\begin{tabular}{lll}
\toprule
Result & Script & Time \\
\midrule
Main ablation, 6 cells, 5 splits & \texttt{run\_multisplit.py} & $\sim$60 min \\
Baseline sanity check & \texttt{run\_baselines.py} & $\sim$5 min \\
Sensitivity curves & \texttt{run\_sensitivity.py} & $\sim$15 min \\
Class F1 and reliability & \texttt{run\_classmetrics.py}, \texttt{run\_reliability.py} & $\sim$4 min \\
Cold-start buckets & \texttt{run\_coldstart.py} & $\sim$2 min \\
Embedding projections & \texttt{run\_visualization.py} & $\sim$1 min \\
MovieLens-1M scale-up & \texttt{run\_ml1m.py} & $\sim$85 min \\
\bottomrule
\end{tabular}
}
\end{table}

\subsection{Discussion}
The ablation paints a sharper picture than the headline number alone. Gated
fusion accounts for most of the empirical RMSE change; the ordinal head
provides native probability estimates cheaply --- four learnable thresholds
--- but does not help RMSE on its own. This interaction suggests that ordinal
heads are most useful here when paired with informative side-information
mechanisms. Limitations are the small benchmark, coarse demographic/genre side
features, and a rating-marginal calibration recipe that does not model
per-user heteroscedasticity.

\section{Conclusion}
\label{sec:conclusion}
We presented a calibration recipe for cumulative-link collaborative
filtering: initialise the monotone thresholds from the training-set rating
quantiles, parameterise their spacing through a softplus-inverse, and
exclude them from $L_2$ regularisation. Paired with per-dimension gated side
information, the resulting model attains RMSE $0.9124\pm0.0036$ and NLL
$1.2584$ across the five MovieLens-100K splits, improving the across-split
mean over both a no-fusion MF analogue and additive fusion. The main lesson is
that gating drives the RMSE gain, while the ordinal head gives calibrated
probabilities and better boundary-class behaviour at negligible parameter
cost.

Practically, this means the ordinal layer should be read as a calibration
interface rather than as the sole source of ranking improvement. It is useful
when the backbone already has enough signal to separate rating boundaries, and
it provides NLL and class probabilities that a sigmoid-MSE head cannot provide
natively. The gated side-information path is the larger accuracy mechanism, but
the cold-start breakdown shows that its gain is not uniform across activity
regimes. A production extension would therefore keep the calibrated ordinal
head, but allow the fusion strategy to vary by user activity instead of using
one global gate behaviour everywhere.

\begingroup
\footnotesize
\setlength{\parskip}{0pt}
\setlength{\itemsep}{0pt}
\bibliographystyle{plain}
\bibliography{refs}
\endgroup

\end{document}
